% TOP_PROBLEM: {"id": 9600004, "title": "Negative association for asymmetric exclusion", "category_id": 19, "category": {"id": 19, "name": "probability", "display_name": "Probability", "description": "Problems involving probability theory, stochastic processes, and random structures.", "slug": "probability", "order_index": 19, "created_at": "2026-07-31T15:26:25.671Z"}, "status": "partially_solved", "problem_number": "AMR-095-0004", "set_id": 13, "statusQualification": "Positive reverse rate q implements the irreducibility assumed in the source introduction; the one-way TASEP boundary case is not needed for that formulation."}
% FIRST_POSTED: 2026-10-01
% ENTRY_KIND: statement_only
% UnsolvedMath ID 9600004; source catalogue code AMR-095-0004
% !TeX program = lualatex
% Definition and sources only; this file is not an LLM solution attempt.
\documentclass[11pt,a4paper]{article}
\usepackage[margin=25mm]{geometry}
\usepackage{fontspec}
\setmainfont{lmroman10-regular.otf}[BoldFont=lmroman10-bold.otf,
 ItalicFont=lmroman10-italic.otf,BoldItalicFont=lmroman10-bolditalic.otf]
\usepackage{amsmath,amssymb,mathtools}
\usepackage{unicode-math}
\setmathfont{latinmodern-math.otf}
\AtBeginDocument{\let\setminus\smallsetminus}
\usepackage{xurl,hyperref}
\hypersetup{colorlinks=true,urlcolor=blue}
\setcounter{secnumdepth}{0}
\setlength{\parindent}{0pt}
\setlength{\parskip}{5pt}
\setlength{\emergencystretch}{3em}
\begin{document}
\section*{Negative association for asymmetric exclusion}\label{9600004}
{\small UnsolvedMath ID 9600004 · AMR-095-0004 · Probability · AMR Open Problem Lists}
\subsection{Definitions and mathematical statement}
Let $p>q>0$ with $p+q=1$. On $\mathbb Z$, run continuous-time nearest-neighbor exclusion: a particle attempts a step to the right at rate $p$ and to the left at rate $q$, and succeeds only if the destination is vacant. The starting configuration is the deterministic step
\[
\eta_0(x)=\begin{cases}1,&x\le0,\\0,&x\ge1.\end{cases}
\]
Let $\mu_t$ be the distribution of the occupation variables $\eta_t\in\{0,1\}^{\mathbb Z}$ at time $t$.

A law $\mu$ is negatively associated if, whenever $A,B\subset\mathbb Z$ are disjoint finite sets and $f:\{0,1\}^A\to\mathbb R$, $g:\{0,1\}^B\to\mathbb R$ are coordinatewise increasing, one has
\[
\mathbb E_\mu[f(\eta|_A)g(\eta|_B)]
\le \mathbb E_\mu[f(\eta|_A)]\,
     \mathbb E_\mu[g(\eta|_B)].
\]
Is $\mu_t$ negatively associated for every $t\ge0$ and every such asymmetric kernel? These finite-coordinate inequalities are the standard definition; extension to bounded increasing functions on disjoint infinite coordinate sets follows by approximation when appropriate.

The orientation of the initial step is part of the problem: particles initially occupy the half-line opposite the direction of drift and spread into empty sites to their right. The question concerns all increasing observables on disjoint sets, so proving nonpositive covariance for two individual sites alone would not establish the assertion. Neither stationarity nor a limit as $t\to\infty$ is assumed; the required inequality is at each fixed time, for the full transient distribution.
\subsection{Short English statement}
Does exclusion begun with all sites on the left occupied preserve negative association while particles drift to the right?
\subsection{Sources}
\begin{itemize}
\item[S1] ULAM.AI and the UnsolvedMath Contributors, supplied problems.json, record 9600004 (AMR-095-0004), AMR Open Problem Lists. Catalogue identity and source transcription; CC BY 4.0.
\url{https://huggingface.co/datasets/ulamai/UnsolvedMath}.
\item[S2] Liggett - Some Open Problems (2012), Problem 4, PDF page 2. Original source indexed by AMR-095-0004.
\url{https://web.archive.org/web/20150919235903id_/http://www.math.ucla.edu/~tml/open.pdf}.
\end{itemize}
\end{document}
